\documentclass[10pt,a4paper]{amsart}
\usepackage[margin=19mm]{geometry}
\usepackage{amsmath,amssymb,mathtools}
\usepackage[scr=rsfs,cal=euler]{mathalfa}
\usepackage{xfrac}
\usepackage[unicode,hidelinks,bookmarks=false]{hyperref}
\newcommand{\ie}{i.e.\ }
\newcommand{\cf}{cf.\ }
\newcommand{\resp}{respectively\ }
\newcommand{\Subsection}{\subsection}
\providecommand{\nobreakdash}{\nobreak\hbox{-}}
\setlength{\emergencystretch}{2em}
\multlinegap0pt
\usepackage{bm}
\usepackage{bbm}
\usepackage{dsfont}
\usepackage{xspace}
\usepackage{cancel}
\usepackage{mathtools}
\usepackage{amsthm}
\makeatletter
\@ifundefined{theorem}{\newtheorem{theorem}{Theorem}}{}
\@ifundefined{lemma}{\newtheorem{lemma}[theorem]{Lemma}}{}
\makeatother
\DeclareMathOperator{\spt}{sp}
\newcommand{\D}{\mathcal{D}}
\newcommand{\Lip}{{\mathrm{lip}}}
\newcommand{\Mes}{{\mathcal{M}}}
\newcommand{\Rd}{\R^d}
\newcommand{\R}{\mathbb{R}}
\newcommand{\Sdd}{{\mathcal{S}^{d \times d}}}
\newcommand{\abs}[1]{{\left \lvert #1 \right \rvert}}
\newcommand{\dive}{{\mathrm{div}}}
\newcommand{\pairing}[1]{{\left \langle #1 \right \rangle}}
\title[Kantorovich-Rubinstein duality theory for the Hessian]{Kantorovich-Rubinstein duality theory for the Hessian}
\author{Karol Bo{\l}botowski, Guy Bouchitt\'e}
\date{Source: arXiv:2412.00516v3 (CC BY 4.0)}
\begin{document}
\maketitle

\noindent\emph{Source and licence.} Kantorovich-Rubinstein duality theory for the Hessian. Version arXiv:2412.00516v3 (CC BY 4.0), distributed under the Creative Commons Attribution 4.0 International licence, as recorded in the arXiv licence metadata for this exact version and shown on the abstract page https://arxiv.org/abs/2412.00516v3. Full rights notice: licenses/arxiv-2412.00516-CC-BY-4.0.txt.\par\smallskip

\noindent\textit{Editorial scope.} Counted unit: the Lemma whose source label is \texttt{lemma:byparts} in the article (source0.tex lines 1933--1940, source label \texttt{lemma:byparts}), delivered together with the article's own complete original proof of it (source0.tex lines 1941--1987, one \texttt{proof} environment of 47 lines). No mathematics is written, repaired or re-derived: statement and proof are the article's own text line for line. Required context retained uncounted: domiCV (lines 736--741); genload (lines 747--750). Every definition, display and label of those lines is retained unchanged. Omitted from this package and not redistributed: 1--735, 742--746, 751--1932, 1988--2383. Nothing inside a retained excerpt is omitted or altered: every retained body is delivered exactly as the article's lines are written, so no omission receipt of any kind accompanies this package. Citations: the retained excerpts carry no citation command at all, so no bibliography entry of the article is transcribed and no reference is redistributed. Reference adaptation: none was needed and none was made; every reference inside the delivered unit resolves inside it, so no unresolved reference or citation is produced. Printed slips kept literal: no printed word, symbol, hypothesis or number of the counted statement or of its proof was altered. Layout and class: the article's own class is \texttt{amsart} with packages including amsmath, amsthm, amssymb, amsfonts, epsfig, color, graphicx, enumerate, enumitem, accents, rotating, multirow, url, mathrsfs; the wrapper is the lane's pinned \texttt{amsart} editorial wrapper at 10pt on A4, which restates the theorem-like environments the retained text needs and adds the article's own macro definitions that the retained text needs (\texttt{\textbackslash D}, \texttt{\textbackslash Lip}, \texttt{\textbackslash Mes}, \texttt{\textbackslash R}, \texttt{\textbackslash Rd}, \texttt{\textbackslash Sdd}, \texttt{\textbackslash abs}, \texttt{\textbackslash dive}, \texttt{\textbackslash pairing}, \texttt{\textbackslash spt}), with extra wrapper packages (bm, bbm, dsfont, xspace, cancel, mathtools). The delivered numbering is the unit's own. Assets: no retained span contains a figure environment, a graphics-inclusion command, a PDF-page inclusion, a table or a TikZ picture; no image file is copied into this package, the package declares no asset dependency and no image file is redistributed with it. Licence: version arXiv:2412.00516v3 of this article is distributed under the Creative Commons Attribution 4.0 International licence, as recorded in the arXiv licence metadata for this exact version and shown on the abstract page https://arxiv.org/abs/2412.00516v3. A full rights notice is delivered at licenses/arxiv-2412.00516-CC-BY-4.0.txt. Characters: every retained excerpt is pure ASCII, so the delivered engine, XeLaTeX with its default Latin Modern text font, reports no missing character.
\begin{align}\label{domiCV}
\sup_n \| v_n\|_{X_p} <+\infty \quad \text{and}\quad \lim_{n\to \infty} & v_n(x) =0 \quad \forall \,x\in\R^d \\
\nonumber
 &\Rightarrow \quad \ \ \pairing{v_n,\mu}  \to 0 \quad \ \ 
 \forall\, \mu\in\Mes_p(\R^d) .
\end{align}

\begin{equation}\label{genload}
  \int f_0=0, \qquad  {\int x f_0 + \int F=0
} .
\end{equation}

\begin{lemma}\label{lemma:byparts}  Let  $f = f_0 - \dive\, F$, where
 $(f_0, F)$ is any pair in $\Mes_2(\R^d)\times\Mes_1(\R^d;\R^d)$ satisfying \eqref{genload}. Let $\sigma \in \Mes(\R^d;\Sdd)$ 
 satisfy  $\dive^2 \sigma = f$ in  $\D'(\R^d)$.
 Then, for every $u\in C^2(\R^d)$ with $\Lip(\nabla u)<+\infty$, we have,
\begin{equation}\label{Rd-byparts}
 { \pairing {\nabla^2 u, \sigma}}  =   \pairing{u,f}  = \pairing{u,f_0} + \pairing{\nabla u, F}.
\end{equation}
\end{lemma}

\begin{proof} By the orthogonality conditions \eqref{genload}, \eqref{Rd-byparts} is valid for affine functions. Therefore, it is not restrictive to assume that $u$ and $\nabla u$ vanish at $0$; we may also assume that $\Lip(\nabla u)\le 1$.
On the other hand, the  equality $\dive^2\sigma=f$ in the sense of distributions implies that  \eqref{Rd-byparts} holds true if $u\in \D(\R^d)$. By using smooth convolution kernels, this can be extended to $u\in C^2(\R^d)$ that are compactly supported in $\R^d$. In order to remove the latter condition,   we consider a sequence of radial  cut-off functions $\eta_k(x) := \eta(\frac{|x|}{k})$ where,
 $$\eta\in \D(\R;[0,1]),\qquad \eta(t)=1 \quad \text{if\quad$|t|\le k$} ,\qquad \eta(t)=0 \quad \text{if\quad$t\ge 2k$}.$$
Then, we set $u_k := u\, \eta_k$.  Since  $u_k $  satisfies  \eqref{Rd-byparts}, we have only to check that the sequence $(v_k)$, given by $v_k := u- u_k= (1-\eta_k) u$,    satisfies,
\begin{equation}\label{claim-truncature}
\pairing{v_k,f_0}\to 0,\qquad \pairing{\nabla v_k,F}\ \to 0, \qquad \pairing{\nabla^2 v_k,\sigma} \to 0 .
\end{equation}
Since $v_k(x)$ vanishes for $|x|\le k$, while $|v_k(x)| \le |(u(x)| \le \frac1{2} |x|^2$  elsewhere, we infer that 
$(v_k)$ is bounded in $X_2(\R^d)$. Hence, $\pairing{v_k,f_0}\to 0 $ by applying \eqref{domiCV} with $\mu=f_0 \in \Mes_2(\R^d)$.
In the same way, $\nabla v_k$ is supported on the subset $\{k\le |x|\le 2k\}$, where it satisfies the upper bound,
\begin{align*}
	|\nabla v_k(x)| &= \left|\big(1-\eta_k(x)\big) \nabla u(x) -  \frac1{k}  u(x)\,  \eta' \Big(\frac{|x|}{k}\Big) \frac{x}{\abs{x}}\right| \\
	&\le |\nabla u(x)| + \frac{\Lip(\eta)}{k} |u(x)|  \le \big(1 +  \Lip(\eta)\big)\, |x| .
\end{align*}
In the last inequality we used the fact that $|\nabla u(x)| \le |x|$, and  $  |u(x)| \le \frac1{2} |x|^2 \le k |x| $
on $\spt (\nabla v_k)$. 
It follows that  $(\nabla v_k)$ is bounded in $X_1(\R^d;\Rd)$ and, recalling that  $F\in \Mes_1(\R^d;\R^d)$, we may apply 
\eqref{domiCV} to infer that $\pairing{\nabla v_k,F} \to 0$.
Next,  we see that\  $\nabla^2 v_k= (1-\eta_k) \nabla^2 u  - (\nabla \eta_k \otimes \nabla u +\nabla u\otimes \nabla \eta_k
+ u  \nabla^2 \eta_k)$ where:
\begin{align*}
	\nabla  \eta_k (x) &=  \frac1{k} \eta'\Big(\frac{|x|}{k}\Big) \frac{x}{|x|},\\
	 \nabla^2 \eta_k(x) &=\frac1{k^2} \eta''\Big(\frac{|x|}{k}\Big) \frac{x\otimes x}{|x|^2} + \frac1{k|x|} \eta' \Big(\frac{|x|}{k}\Big)\left( \mathrm{Id} - \frac{x\otimes x}{|x|^2}\right).
\end{align*}
 We notice  that:
 \begin{enumerate}
\item[-] for any $x \neq 0$ there hold the bounds $\varrho(\nabla \eta_k \otimes \nabla u +\nabla u\otimes \nabla \eta_k)(x) \le 2|\nabla u(x)|\frac{\Lip(\eta)}{k}$ and $\varrho\big(\nabla^2 \eta_k(x)\big) \leq \frac{\Lip(\eta')}{k^2} + \frac{\Lip(\eta)}{k \abs{x}}$;


\item[-] $\nabla \eta_k$ and $\nabla^2 \eta_k$ are supported on $\{ k\le |x|\le 2k\}$, where it holds that $|u(x)|\le \frac1{2} |x|^2\le 2 k^2$ and $|\nabla u(x)|\le  |x|\le 2 k$.
\end{enumerate}
All in all,  we obtain a uniform upper bound for $\varrho(\nabla^2 v_k)$ whose support is contained  in $\{k\le |x| \le 2k\}$,
\begin{align*} \varrho\big(\nabla^2 v_k(x)\big) &\le  \varrho\big(\nabla^2u(x)\big)  + \varrho\big(\nabla \eta_k \otimes \nabla u +\nabla u\otimes \nabla \eta_k\big)(x)
+ |u(x)|\, \varrho\big(\nabla^2 \eta_k(x)\big) \\
&\le 1 +  2\,|\nabla u(x)|\frac{\Lip(\eta)}{k} + |u(x)| \Big( \frac{\Lip(\eta')}{k^2} + \frac{\Lip(\eta)}{k \abs{x}} \Big)\\
& \le 1 +  2 \,\Lip(\eta) \frac{|x|}{k} + \frac1{2} \Lip(\eta') \frac {|x|^2}{k^2} + \frac1{2} \Lip(\eta) \frac {|x|}{k} \\
&\le C:=1 + 5 \, \Lip(\eta) + 2 \,\Lip(\eta').
\end{align*} 
By virtue of the  inequality  $|\pairing{\nabla^2 v_k,\sigma}|\, \le\,  \varrho(\nabla^2 v_k)\, \varrho^0(\sigma) \le C  \varrho^0(\sigma) $ 
holding in  the sense of measures,
it follows that,
$$  |\pairing{\nabla^2 v_k,\sigma}|  \le  C \int_{\{k\le |x| \le 2k\}}  \varrho^0(\sigma) .
$$
Since $\int_{\R^d} \varrho^0(\sigma)<+\infty$, we conclude that $\pairing{\nabla^2 v_k,\sigma}\to 0$ for $k\to\infty$
as required in \eqref{claim-truncature}.
This ends the proof.
\end{proof}

\end{document}
